\documentclass[aspectratio=169]{beamer}
\usepackage{amsmath,amssymb}
\usepackage{graphicx}
\usepackage{booktabs}
\usepackage{hyperref}

\usetheme{Madrid}
\usecolortheme{default}

\title{ACP2 - Stage 3: Set Valued Dynamical System Visualization}
\author{Nam Do, Jyri Ollanketo\\Supervised by: Dr. Kalle Timperi}
\institute{University of Oulu}
\date{April 14, 2026}

\begin{document}

\begin{frame}
  \titlepage
\end{frame}

\begin{frame}{Implemented features in stage 3}
  \begin{itemize}
    \item Continuous dynamical system visualization
    \item Extensibility with arbitrary equations
    \item GAIO based Ulam method for continuous time
    \item Parameter sweeping for discrete time
    \item GUI update and usability improvements
  \end{itemize}
\end{frame}

\section{Continuous Visualization}

\begin{frame}{Continuous Dynamical System Visualization}
  \textbf{Goal:} extend set valued visualization from discrete maps to continuous flows.

  \vspace{6pt}
  \textbf{Duffing ODE:}
  \begin{equation}
    \dot{x} = y, \quad \dot{y} = x - x^3 - \delta y.
  \end{equation}

  \vspace{4pt}
  \textbf{Boundary differential equation:}
  \begin{align}
    \dot{x} &= f(x) + \epsilon n, \\
    \dot{n} &= -Df(x)^T n + \bigl(Df(x)^T n \cdot n\bigr) n.
  \end{align}

  \vspace{4pt}
  \begin{itemize}
    \item RK4 integration for $(x,n)$ and arc length reparameterization
    \item Visualization of evolving boundary curves in the phase plane
  \end{itemize}
\end{frame}

\section{Extensibility}

\begin{frame}{Extensibility for Arbitrary Simple Equations}
  \textbf{Feature:} user defined systems for both discrete and continuous cases.

  \vspace{4pt}
  \begin{itemize}
    \item Two equation parser with parameter list validation
    \item Reserved identifiers blocked for safety and clarity
    \item Same backend pipeline for custom and built in systems
  \end{itemize}

  \vspace{6pt}
  \textbf{Example custom system:}
  \begin{equation}
    x' = 1 - a x^2 + y, \quad y' = b x.
  \end{equation}

  \textbf{Example custom ODE:}
  \begin{equation}
    \dot{x} = y, \quad \dot{y} = x - x^3 - \delta y.
  \end{equation}
\end{frame}

\section{Continuous Ulam}

\begin{frame}{GAIO Based Ulam Method for Continuous Time}
  \textbf{Goal:} approximate invariant measures for flows.

  \vspace{4pt}
  \begin{itemize}
    \item Partition domain into grid boxes $\{B_i\}$
    \item Build transition probabilities via time $T$ flow map
    \item Use epsilon inflation for set valued images
  \end{itemize}

  \vspace{6pt}
  \begin{equation}
    P_{ij} = \frac{1}{N_i} \#\{x_k \in B_i : B_{\epsilon}(\varphi_T(x_k)) \cap B_j \ne \emptyset\}.
  \end{equation}

  \vspace{4pt}
  \begin{itemize}
    \item Right eigenvector gives invariant measure
    \item Left eigenvector highlights dual repeller structure
  \end{itemize}
\end{frame}

\section{Parameter Sweep}

\begin{frame}{Parameter Sweeping for Discrete Time}
  \textbf{Goal:} explore bifurcations and stability changes efficiently.

  \vspace{4pt}
  \begin{itemize}
    \item Sweep a parameter over a range
    \item Compute periodic orbits and classify stability
    \item Export CSV and JSON for post analysis
  \end{itemize}

  \vspace{6pt}
  \textbf{Outputs:}
  \begin{itemize}
    \item Per sample counts of stable, saddle, and unstable orbits
    \item Orbit locations in extended boundary space
  \end{itemize}
\end{frame}

\section{GUI Update}

\begin{frame}{GUI Update and Usability Improvements}
  \textbf{Main updates:}
  \begin{itemize}
    \item System type selector: discrete versus continuous
    \item Axis range control with clamped normalization
    \item Custom equation editor and parameter panel
    \item Parameter sweep panel with summaries and export
    \item Ulam overlay controls and color scale
  \end{itemize}
\end{frame}

\section{Summary}

\begin{frame}{Summary}
  \begin{itemize}
    \item Continuous time boundary visualization is now supported
    \item Arbitrary simple equations enable quick model extension
    \item Continuous Ulam method provides invariant measure overlays
    \item Parameter sweeping speeds up bifurcation exploration
    \item GUI updates improve clarity and workflow
  \end{itemize}
\end{frame}

\begin{frame}{References}
  \begin{thebibliography}{9}
    \bibitem{tey2023} W. H. Tey, \textit{Boundary mappings of minimal invariant sets}, Ph.D. dissertation, Imperial College London, 2023.
    \bibitem{dellnitz2001} M. Dellnitz, G. Froyland, and O. Junge, \textit{The algorithms behind GAIO, set oriented numerical methods for dynamical systems}, Springer, 2001.
    \bibitem{ulam1960} S. M. Ulam, \textit{A Collection of Mathematical Problems}, Interscience Publishers, 1960.
  \end{thebibliography}
\end{frame}

\end{document}
